\documentclass[10pt,a4paper]{amsart}
\usepackage{amsmath,amssymb,amsthm,graphicx}
% Retained minimal preamble for the Stacks Project chapter varieties.tex.
% Verbatim excerpt of the chapter preamble: Xy-pic setup (native preamble lines 14-18), the native theorem
% declarations with the mathematical macros (native preamble lines 154-217), and the project own no-op
% definitions of the reference/slogan/history display environments (native wrapper.tex lines 1-3) that the
% counted statement uses for its slogan. Copied mechanically; nothing added, removed or redefined.
% For commutative diagrams we use Xy-pic
\usepackage[all]{xy}

% We use 2cell for 2-commutative diagrams.
\xyoption{2cell}
% Theorem environments.
%
\theoremstyle{plain}
\newtheorem{theorem}[subsection]{Theorem}
\newtheorem{proposition}[subsection]{Proposition}
\newtheorem{lemma}[subsection]{Lemma}

\theoremstyle{definition}
\newtheorem{definition}[subsection]{Definition}
\newtheorem{example}[subsection]{Example}
\newtheorem{exercise}[subsection]{Exercise}
\newtheorem{situation}[subsection]{Situation}

\theoremstyle{remark}
\newtheorem{remark}[subsection]{Remark}
\newtheorem{remarks}[subsection]{Remarks}

\numberwithin{equation}{subsection}

% Macros
%
\def\lim{\mathop{\mathrm{lim}}\nolimits}
\def\colim{\mathop{\mathrm{colim}}\nolimits}
\def\Spec{\mathop{\mathrm{Spec}}}
\def\Hom{\mathop{\mathrm{Hom}}\nolimits}
\def\Ext{\mathop{\mathrm{Ext}}\nolimits}
\def\SheafHom{\mathop{\mathcal{H}\!\mathit{om}}\nolimits}
\def\SheafExt{\mathop{\mathcal{E}\!\mathit{xt}}\nolimits}
\def\Sch{\mathit{Sch}}
\def\Mor{\mathop{\mathrm{Mor}}\nolimits}
\def\Ob{\mathop{\mathrm{Ob}}\nolimits}
\def\Sh{\mathop{\mathit{Sh}}\nolimits}
\def\NL{\mathop{N\!L}\nolimits}
\def\CH{\mathop{\mathrm{CH}}\nolimits}
\def\proetale{{pro\text{-}\acute{e}tale}}
\def\etale{{\acute{e}tale}}
\def\QCoh{\mathit{QCoh}}
\def\Ker{\mathop{\mathrm{Ker}}}
\def\Im{\mathop{\mathrm{Im}}}
\def\Coker{\mathop{\mathrm{Coker}}}
\def\Coim{\mathop{\mathrm{Coim}}}

% Boxtimes
%
\DeclareMathSymbol{\boxtimes}{\mathbin}{AMSa}{"02}

%
% Macros for moduli stacks/spaces
%
\def\QCohstack{\mathcal{QC}\!\mathit{oh}}
\def\Cohstack{\mathcal{C}\!\mathit{oh}}
\def\Spacesstack{\mathcal{S}\!\mathit{paces}}
\def\Quotfunctor{\mathrm{Quot}}
\def\Hilbfunctor{\mathrm{Hilb}}
\def\Curvesstack{\mathcal{C}\!\mathit{urves}}
\def\Polarizedstack{\mathcal{P}\!\mathit{olarized}}
\def\Complexesstack{\mathcal{C}\!\mathit{omplexes}}
% \Pic is the operator that assigns to X its picard group, usage \Pic(X)
% \Picardstack_{X/B} denotes the Picard stack of X over B
% \Picardfunctor_{X/B} denotes the Picard functor of X over B
\def\Pic{\mathop{\mathrm{Pic}}\nolimits}
\def\Picardstack{\mathcal{P}\!\mathit{ic}}
\def\Picardfunctor{\mathrm{Pic}}
\def\Deformationcategory{\mathcal{D}\!\mathit{ef}}
\newenvironment{reference}{}{}
\newenvironment{slogan}{}{}
\newenvironment{history}{}{}
\title{An automorphism of a reduced finite-type scheme that induces the identity on the underlying topological space acts trivially on coherent modules and is the identity in positive dimension}
\author{The Stacks Project Authors}
\date{stacks-project:varieties.tex:local-snapshot-20261011; GNU FDL 1.2 or later}
\begin{document}
\maketitle
\noindent\textit{Editorial note.} All delivered mathematical text is the author's own, copied verbatim from the pinned Stacks Project chapter file varieties.tex; no statement or proof step was rewritten, reordered or removed.\par\medskip
\noindent\textit{Editorial note.} The proof cites no external \textbackslash{}ref inside the counted span, so the assembly needs no external-reference stubs.\par\medskip
\noindent\textit{Editorial note.} The author's parenthetical 'details omitted' remarks concern only auxiliary facts (the span of constant functions and finiteness of the automorphism group of a curve over a finite field); the main case arguments are fully present.\par\medskip
\noindent\textit{Allowed standard prerequisites (editorial).}\begin{itemize}
\item The set of closed points and generic points of a scheme of finite type over a field. Stacks Project, Properties of Schemes and Varieties (topology of schemes of finite type over a field). Background for the reductions to irreducible components and to dense affine opens made in the proof.
\item A curve over a finite field has a finite automorphism group. Stacks Project, Varieties (finiteness of automorphism groups of curves over finite fields). Invoked in the finite-field case to give f finite order, on which the prime-to-order density argument depends (proof line 6131).
\end{itemize}\medskip
\begin{lemma}
\label{lemma-automorphism}
Let $X$ be a reduced scheme of finite type over a field $k$. Let $f : X \to X$
be an automorphism over $k$ which induces the identity map on the underlying
topological space of $X$. Then
\begin{enumerate}
\item $f^*\mathcal{F} \cong \mathcal{F}$ for every coherent
$\mathcal{O}_X$-module, and
\item if $\dim(Z) > 0$ for every irreducible component $Z \subset X$,
then $f$ is the identity.
\end{enumerate}
\end{lemma}

\begin{proof}
Part (1) follows from part (2) and the fact that the connected components
of $X$ of dimension $0$ are spectra of fields.

\medskip\noindent
Let $Z \subset X$ be an irreducible component viewed as an integral closed
subscheme. Clearly $f(Z) \subset Z$ and $f|_Z : Z \to Z$ is an automorphism
over $k$ which induces the identity map on the underlying topological space
of $Z$. Since $X$ is reduced, it suffices to show that the arrows
$f|_Z : Z \to Z$ are the identity. This reduces us to the case discussed
in the next paragraph.

\medskip\noindent
Assume $X$ is irreducible of dimension $> 0$. Choose a nonempty
affine open $U \subset X$. Since $f(U) \subset U$ and since
$U \subset X$ is scheme theoretically dense it suffices to
prove that $f|_U : U \to U$ is the identity.

\medskip\noindent
Assume $X = \Spec(A)$ is affine, irreducible, of dimension $> 0$
and $k$ is an infinite field. Let $g \in A$ be nonconstant. The set
$$
S = \bigcup\nolimits_{\lambda \in k} V(g - \lambda)
$$
is dense in $X$ because it is the inverse image of the dense subset
$\mathbf{A}^1_k(k)$ by the nonconstant morphism $g : X \to \mathbf{A}^1_k$.
If $x \in S$, then the image $g(x)$ of $g$ in $\kappa(x)$
is in the image of $k \to \kappa(x)$. Hence
$f^\sharp : \kappa(x) \to \kappa(x)$ fixes $g(x)$.
Thus the image of $f^\sharp(g)$ in $\kappa(x)$ is equal to $g(x)$.
We conclude that
$$
S \subset V(g - f^\sharp(g))
$$
and since $X$ is reduced and $S$ is dense we conclude $g=f^\sharp(g)$.
This proves $f^\sharp = \text{id}_A$ as $A$ is generated as a $k$-algebra
by elements $g$ as above (details omitted; hint: the set of constant
functions is a finite dimensional $k$-subvector space of $A$).
We conclude that $f = \text{id}_X$.

\medskip\noindent
Assume $X = \Spec(A)$ is affine, irreducible, of dimension $> 0$
and $k$ is a finite field. If for every $1$-dimensional integral
closed subscheme $C \subset X$ the restriction $f|_C : C \to C$
is the identity, then $f$ is the identity. This reduces us to the
case where $X$ is a curve. A curve over a finite field has a
finite automorphism group (details omitted). Hence $f$ has finite
order, say $n$. Then we pick $g : X \to \mathbf{A}^1_k$
nonconstant as above and we consider
$$
S = \{x \in X\text{ closed such that }[\kappa(g(x)) : k]
\text{ is prime to }n\}
$$
Arguing as before we find that $S$ is dense in $X$. Since
for $x \in X$ closed the map
$f^\sharp : \kappa(x) \to \kappa(x)$ is an
automorphism of order dividing $n$
we see that for $x \in S$ this automorphism
acts trivially on the subfield generated by
the image of $g$ in $\kappa(x)$. Thus we conclude that
$S \subset V(g - f^\sharp(g))$ and we win as before.
\end{proof}

\end{document}
